\scanpage{62}
\mspage{5}

les propriétés de régularité de $C$ en tant que cercle \inserted{``intérieur''} dans la surface \deleted{\unsure{sans bord}} \inserted{conforme} $X'$ déduite par bouchage du trou. D'autre part, le choix d'un point de $D\smallsetminus C$ équivaut à celui d'une structure torique sur $C$ compatible avec la structure de droite projective réelle (i.e. une restriction du groupe structural \inserted{$\unsure{\mathrm{GP}(1,\mathbb R)|\mathrm{Gp}(1,\mathbb R)}=$} $\mathrm{Gl}(2,\mathbb R)/\mathbb R^{*}$ au sous-groupe cpct maximal $\mathrm O(2)/\pm1$).
--- et cel équivaut aussi à la donnée d'une structure \deleted{\unsure{do}}torique de $D$ sur \inserted{\unsure{la base|le bord}} $\partial D=C$. Notons que l'on aura alors une inclusion canonique
\[
D\subset T_{X',a}
\]
où $a$ est le centre choisi de $D$. On va expliciter \deleted{ainsi} ces dernières relations ainsi

\uline{Prop.} Équivalence entre \deleted{la cat. des surfaces \unsure{confo|conf.}}
% archive sheet number: 64
